\documentclass[14pt]{extarticle}
\usepackage{fontspec}
\usepackage{polyglossia}
\usepackage{amsmath}
\usepackage{amssymb}
\usepackage{xcolor}
\usepackage{fancyhdr}
\usepackage{graphicx}
\usepackage{listings}
\usepackage[most]{tcolorbox}
% \usepackage{setspace}
\usepackage{pifont}
\usepackage{enumitem}
\usepackage{titlesec}
\usepackage[bottom]{footmisc}
\usepackage{titling}
\usepackage{minted}
\usepackage{etoolbox}
\usepackage{array}
\usepackage{longtable} % page-breakable tables (must load before polyglossia/bidi RTL setup)
\usepackage{extsizes}
\usepackage{float}
\usepackage{booktabs}
\usepackage{hyperref}
\usepackage[onehalfspacing]{setspace}
\newif\ifwhiteblack
\ifdefined\setwhiteblack
  \setwhiteblack
\fi

\ifwhiteblack
  \usepackage[a4paper, landscape, top=1.25cm, bottom=1.25cm, left=1.25cm, right=1.25cm, headheight=1.5cm, headsep=0.5cm]{geometry}
\else
  \usepackage[a4paper, top=2.7cm, bottom=2.7cm, left=2.5cm, right=2.5cm, headheight=1.5cm, headsep=0.5cm]{geometry}
\fi
\usepackage{tikz}
\usepackage{datetime2}

\usetikzlibrary{automata, positioning, arrows.meta, arrows}
\providecommand{\LTcaptype}{table}

\newfontfamily\emoji{Segoe UI Emoji}

\pagestyle{fancy}

\setmainlanguage[numerals=western]{arabic}
\setotherlanguage{english}
\newfontfamily\arabicfont[Script=Arabic]{Calibri}
\newfontfamily\arabicfonttt[Script=Arabic]{Courier New}



\ifwhiteblack
\lstset{
  language=[Sharp]C,
  numbers=left,
  stepnumber=1,
  numbersep=8pt,
  frame=single,
  basicstyle=\ttfamily\small,
  keywordstyle=\bfseries,
  stringstyle=\ttfamily,
  commentstyle=\itshape
}
\else
\lstset{
  language=[Sharp]C,
  numbers=left,
  stepnumber=1,
  numbersep=8pt,
  frame=single,
  basicstyle=\ttfamily\small,
  keywordstyle=\color{blue},
  stringstyle=\color{red},
  commentstyle=\color{green!50!black}
}
\fi

\newif\ifdetailed
\ifdefined\setdetailed
  \setdetailed
\fi

\newif\ifwithsols
\ifdefined\setwithsols
  \setwithsols
\fi

\ifwhiteblack
  \AtBeginDocument{
    \LARGE
    \usemintedstyle{bw}
  }
\fi

% Default minted look: compact code that wraps long lines.
% Not applied in black-and-white landscape mode, which already uses \LARGE.
\ifwhiteblack\else
  \setminted{fontsize=\small, breaklines=true}
\fi

% unified tcolorboxes for articles
\ifwhiteblack
\tcbset{colback=white, colframe=black, fonttitle=\bfseries, boxrule=1pt}
\newtcolorbox{boxDef}[1][]{title={{\emoji📘} تعريف\ifx\\#1\\\else ~#1\fi :}}
\newtcolorbox{boxExercise}[1][]{title={{\emoji🧩} تمرين\ifx\\#1\\\else ~#1\fi :}}
\newtcolorbox{boxExample}[1][]{title={{\emoji📝} مثال\ifx\\#1\\\else ~#1\fi :}}
\newtcolorbox{boxNote}[1][]{title={{\emoji✨} ملاحظة\ifx\\#1\\\else ~#1\fi :}}
\newtcolorbox{boxAttention}[1][]{title={{\emoji🔔} تنبيه\ifx\\#1\\\else ~#1\fi :}}
\newtcolorbox{boxWarning}[1][]{title={{\emoji⚡} ملاحظة هامة\ifx\\#1\\\else ~#1\fi :}}
\newtcolorbox{boxSolution}[1][]{title={{\emoji✅} حل\ifx\\#1\\\else ~#1\fi :}}
\newtcolorbox{boxSymbol}[1][]{title={{\emoji🔣} رمز\ifx\\#1\\\else ~#1\fi :}}
\newtcolorbox{boxHint}[1][]{title={{\emoji💡} تلميح\ifx\\#1\\\else ~#1\fi :}}
\else
\tcbset{colback=white, colframe=black, fonttitle=\bfseries, boxrule=0.8pt}
\newtcolorbox{boxDef}[1][]{colback=blue!5!white,colframe=blue!75!black,
  title={{\emoji📘} تعريف\ifx\\#1\\\else ~#1\fi :}}
\newtcolorbox{boxExercise}[1][]{colback=cyan!5!white,colframe=cyan!70!black,
  title={{\emoji🧩} تمرين\ifx\\#1\\\else ~#1\fi :}}
\newtcolorbox{boxExample}[1][]{colback=yellow!5!white,colframe=orange!90!black,
  title={{\emoji📝} مثال\ifx\\#1\\\else ~#1\fi :}}
\newtcolorbox{boxNote}[1][]{colback=gray!10!white,colframe=black,
  title={{\emoji✨} ملاحظة\ifx\\#1\\\else ~#1\fi :}}
\newtcolorbox{boxAttention}[1][]{colback=magenta!10!white,colframe=magenta!80!black,
  title={{\emoji🔔} تنبيه\ifx\\#1\\\else ~#1\fi :}}
\newtcolorbox{boxWarning}[1][]{colback=red!5!white,colframe=red!75!black,
  title={{\emoji⚡} ملاحظة هامة\ifx\\#1\\\else ~#1\fi :}}
\newtcolorbox{boxSolution}[1][]{colback=green!5!white,colframe=green!60!black,
  title={{\emoji✅} حل\ifx\\#1\\\else ~#1\fi :}}
\newtcolorbox{boxSymbol}[1][]{colback=purple!5!white,colframe=purple!70!black,
  title={{\emoji🔣} رمز\ifx\\#1\\\else ~#1\fi :}}
\newtcolorbox{boxHint}[1][]{colback=teal!5!white,colframe=teal!60!black,
  title={{\emoji💡} تلميح\ifx\\#1\\\else ~#1\fi :}}
\fi


\ifwhiteblack
\tcbset{simplecode/.style={ colback=white, colframe=black, boxrule=1pt, arc=2pt, left=4pt,right=4pt,top=4pt,bottom=4pt}}
\else
\tcbset{simplecode/.style={ colback=gray!5, colframe=black!50, boxrule=0.4pt, arc=2pt, left=4pt,right=4pt,top=4pt,bottom=4pt}}
\fi
% Compact code boxes: tighter line spacing (the preamble uses onehalfspacing) and small vertical gaps.
% Must come AFTER the simplecode definition above.
\tcbset{simplecode/.append style={fontupper=\setstretch{0.95}, top=4pt, bottom=4pt, before skip=6pt, after skip=6pt}}
\newenvironment{boxCode}{\begin{tcolorbox}[simplecode]}{\end{tcolorbox}}

% Code cell used inside tables (keeps code left-to-right inside RTL tables).
% Use inside a fixed-width p{...} column: \begin{pycell}\begin{minted}...\end{minted}\end{pycell}
\newtcolorbox{pycell}{frame empty,colback=white,boxsep=0pt,left=2pt,right=2pt,top=2pt,bottom=2pt,width=\linewidth}

\newcolumntype{C}[1]{>{\centering\arraybackslash}p{#1}}

\makeatletter
\DTMnewdatestyle{dotted}{%
  \renewcommand{\DTMdisplaydate}[4]{\two@digits{##3}.\two@digits{##2}.##1}%
  \renewcommand{\DTMDisplaydate}{\DTMdisplaydate}%
}
\makeatother
\DTMsetdatestyle{dotted}

% redefine spaces after titles
\makeatletter
\renewcommand{\@maketitle}{%
  \begin{center}
    {\huge \bfseries \@title \par}%
    \vskip 0.2em % space between title and author
    {\large \@author \par}%
    % \vskip 0.2em % space between author and date
    % {\normalsize \@date \par}%
    {\small \hfill \begin{english}\DTMtoday\end{english} \par}%
  \end{center}
}
\makeatother

\fancyhf{} % clear default
\fancypagestyle{plain}{
  \fancyhf{}
  \fancyhead[L]{مدرسة التسامح الشاملة}
  % \fancyhead[L]{\includegraphics[height=1cm]{../../../images/logoTasamoh.png}}
  \fancyhead[R]{الأستاذ محمود اغبارية}
  \fancyfoot[C]{\thepage}
}

\fancyhead[L]{مدرسة التسامح الشاملة}
\fancyhead[R]{الأستاذ محمود اغبارية}
\fancyfoot[C]{\thepage}

\setcounter{tocdepth}{3} % only section subsection and subsubsection in TOC

\newcommand{\importpdfpage}[2]{%
    \thispagestyle{empty}
    \newgeometry{left=0cm,right=0cm,top=0cm,bottom=0cm}
    \includegraphics[page=#2,width=0.95\paperwidth,height=0.95\paperheight,keepaspectratio]{#1}
    \restoregeometry
}

\newcommand{\insertFullImg}[1]{%
    \noindent
    \makebox[\textwidth][c]{
        \includegraphics[width=0.9\paperwidth,keepaspectratio]{#1}%
    }%
}%


\ifwithsols
\title{حل ورقة تحضير للامتحان النهائي - المتغيرات، العوامل والشروط (بايثون أ)}
\else
\title{ورقة تحضير للامتحان النهائي - المتغيرات، العوامل والشروط (بايثون أ)}
\fi

\begin{document}

\maketitle
\thispagestyle{fancy}

\section*{القسم الأول: المتغيرات والشروط - أسئلة الدرس}

\begin{enumerate}[label=\textbf{تمرين \arabic*:}, wide=0pt, leftmargin=*, itemsep=2.2em]

% ---------------------------------------------------------------
\item \textbf{تمرين متابعة:} اكتبوا ما هو مُخرج مقطع البرنامج التالي دون تشغيله في برمجية بايثون. يجب الاستعانة بجدول متابعة.

\begin{boxCode}
\begin{english}
\begin{minted}[linenos]{python}
num1 = 5
num2 = num1 + 1
num1 = num2
print("num1 = ", num1, " num2 = ", num2)
\end{minted}
\end{english}
\end{boxCode}

\ifwithsols
\begin{boxSolution}
\begin{center}
\begin{english}
\renewcommand{\arraystretch}{1.4}
\begin{tabular}{|C{1.5cm}|C{2cm}|C{2cm}|C{6cm}|}
\hline
\textarabic{\textbf{السطر}} & \textbf{num1} & \textbf{num2} & \textarabic{\textbf{المخرجات}} \\
\hline
1 & 5 & & \\ \hline
2 & 5 & 6 & \\ \hline
3 & 6 & 6 & \\ \hline
4 & 6 & 6 & \texttt{num1 =\ \ 6\ \ num2 =\ \ 6} \\ \hline
\end{tabular}
\end{english}
\end{center}
\end{boxSolution}
\else
\begin{center}
\begin{english}
\renewcommand{\arraystretch}{1.4}
\begin{tabular}{|C{1.5cm}|C{2cm}|C{2cm}|C{6cm}|}
\hline
\textarabic{\textbf{السطر}} & \textbf{num1} & \textbf{num2} & \textarabic{\textbf{المخرجات}} \\
\hline
 & & & \\ \hline
 & & & \\ \hline
 & & & \\ \hline
 & & & \\ \hline
\end{tabular}
\end{english}
\end{center}
\fi
\filbreak

% ---------------------------------------------------------------
\item \textbf{تمرين متابعة:} اكتبوا ما هو مُخرج مقطع البرنامج التالي دون تشغيله في برمجية بايثون. يجب الاستعانة بجدول متابعة.

\begin{boxCode}
\begin{english}
\begin{minted}[linenos]{python}
num1 = 3
num2 = 10
num1 = num2
num2 = num1
print("num1 = ", num1, "num2 = ", num2)
\end{minted}
\end{english}
\end{boxCode}

\ifwithsols
\begin{boxSolution}
\begin{center}
\begin{english}
\renewcommand{\arraystretch}{1.4}
\begin{tabular}{|C{1.5cm}|C{2cm}|C{2cm}|C{6cm}|}
\hline
\textarabic{\textbf{السطر}} & \textbf{num1} & \textbf{num2} & \textarabic{\textbf{المخرجات}} \\
\hline
1 & 3 & & \\ \hline
2 & 3 & 10 & \\ \hline
3 & 10 & 10 & \\ \hline
4 & 10 & 10 & \\ \hline
5 & 10 & 10 & \texttt{num1 =\ \ 10 num2 =\ \ 10} \\ \hline
\end{tabular}
\end{english}
\end{center}
\end{boxSolution}
\else
\begin{center}
\begin{english}
\renewcommand{\arraystretch}{1.4}
\begin{tabular}{|C{1.5cm}|C{2cm}|C{2cm}|C{6cm}|}
\hline
\textarabic{\textbf{السطر}} & \textbf{num1} & \textbf{num2} & \textarabic{\textbf{المخرجات}} \\
\hline
 & & & \\ \hline
 & & & \\ \hline
 & & & \\ \hline
 & & & \\ \hline
 & & & \\ \hline
\end{tabular}
\end{english}
\end{center}
\fi
\filbreak

% ---------------------------------------------------------------
\item \textbf{تمرين متابعة:} اكتبوا ما هو مُخرج مقطع البرنامج التالي دون تشغيله في برمجية بايثون. يجب الاستعانة بجدول متابعة.

\begin{boxCode}
\begin{english}
\begin{minted}[linenos]{python}
counter = 0
print(counter == 5)
counter = counter + 1
print("Hello!")
\end{minted}
\end{english}
\end{boxCode}

\ifwithsols
\begin{boxSolution}
\begin{center}
\begin{english}
\renewcommand{\arraystretch}{1.4}
\begin{tabular}{|C{1.5cm}|C{2.5cm}|C{4cm}|}
\hline
\textarabic{\textbf{السطر}} & \textbf{counter} & \textarabic{\textbf{المخرجات}} \\
\hline
1 & 0 & \\ \hline
2 & 0 & \texttt{False} \\ \hline
3 & 1 & \\ \hline
4 & 1 & \texttt{Hello!} \\ \hline
\end{tabular}
\end{english}
\end{center}
\end{boxSolution}
\else
\begin{center}
\begin{english}
\renewcommand{\arraystretch}{1.4}
\begin{tabular}{|C{1.5cm}|C{2.5cm}|C{4cm}|}
\hline
\textarabic{\textbf{السطر}} & \textbf{counter} & \textarabic{\textbf{المخرجات}} \\
\hline
 & & \\ \hline
 & & \\ \hline
 & & \\ \hline
 & & \\ \hline
\end{tabular}
\end{english}
\end{center}
\fi
\filbreak

% ---------------------------------------------------------------
\item \textbf{تمرين متابعة:} أمامكم مقطع البرنامج التالي:

\begin{boxCode}
\begin{english}
\begin{minted}[linenos]{python}
number_int = int(input("Enter an integer: "))
number_float = float(number_int)
print("The decimal number is:", number_float)
\end{minted}
\end{english}
\end{boxCode}

\begin{enumerate}[label=\alph*.]
\item اكتبوا ما هو مُخرج مقطع البرنامج دون تشغيله في برمجية بايثون، على فرض أنّ المُدخل هو $34$. يجب الاستعانة بجدول متابعة.

\ifwithsols
\begin{boxSolution}
\begin{center}
\begin{english}
\renewcommand{\arraystretch}{1.4}
\begin{tabular}{|C{1.2cm}|C{2.4cm}|C{2.8cm}|C{3.8cm}|}
\hline
\textarabic{\textbf{السطر}} & \textbf{\small number\_int} & \textbf{\small number\_float} & \textarabic{\textbf{المخرجات}} \\
\hline
1 & 34 & & \small\texttt{Enter an integer:} \\ \hline
2 & 34 & 34.0 & \\ \hline
3 & 34 & 34.0 & \small\texttt{The decimal number is: 34.0} \\ \hline
\end{tabular}
\end{english}
\end{center}
\end{boxSolution}
\else
\begin{center}
\begin{english}
\renewcommand{\arraystretch}{1.4}
\begin{tabular}{|C{1.2cm}|C{2.4cm}|C{2.8cm}|C{3.8cm}|}
\hline
\textarabic{\textbf{السطر}} & \textbf{\small number\_int} & \textbf{\small number\_float} & \textarabic{\textbf{المخرجات}} \\
\hline
 & & & \\ \hline
 & & & \\ \hline
 & & & \\ \hline
\end{tabular}
\end{english}
\end{center}
\fi
\filbreak

\item ما الهدف من مقطع البرنامج؟

\ifwithsols
\begin{boxSolution}
يستقبل مقطع البرنامج عددًا صحيحًا من المستخدم، يحوّله من نمط عدد صحيح (\textenglish{int}) إلى نمط عدد عشري (\textenglish{float})، ثم يطبع العدد العشري.
\end{boxSolution}
\fi
\filbreak
\end{enumerate}
\filbreak

% ---------------------------------------------------------------
\item في دكان \textbf{"مفاجآت"}، قرّروا منح هدية لكل زبون وفقًا لمبلغ الشراء من الدكان. الهدية هي واحدة من الخيارات الخمسة التالية:
دفتر يوميات (\textenglish{Diary})، مقلمة مجهّزة (\textenglish{Pencil case})، حقيبة مجهّزة (\textenglish{Bag})، لعبة حاسوب (\textenglish{Computer game})، قسيمة شراء بقيمة $500$ شيكل (\textenglish{Gift voucher}).
عند الشراء بأقل من $10$ شيكل، فإنّ الزبون \textbf{لا يستحق} هدية.

يتم تحديد الهدية وفقًا للسيناريو الذي أمامكم (المتغير \textenglish{amount} هو مبلغ الشراء):

\begin{boxCode}
\begin{english}
\begin{minted}{python}
if amount < 10:
    print("Not eligible for a gift")
elif amount < 400:
    if amount >= 10 and amount <= 100:
        print("Diary")
    elif amount < 250:
        print("Pencil case")
    else:
        print("Bag")
elif amount <= 600:
    print("Computer game")
else:
    print("Gift voucher")
\end{minted}
\end{english}
\end{boxCode}

\begin{enumerate}[label=\alph*.]
\item قرّر يوسف أن يشغّل السيناريو. ساعدوه وأكملوا الجدول أدناه: لكل مبلغ شراء، ما هي الهدية التي يستحقها الزبون؟

{\centering
\renewcommand{\arraystretch}{1.4}
\begin{tabular}{|C{3cm}|C{6cm}|}
\hline
\textbf{مبلغ الشراء} & \textbf{الهدية التي يستحقها الزبون} \\
\hline
\ifwithsols
$8$ & لا يستحق هدية \\ \hline
$80$ & دفتر يوميات \\ \hline
$250$ & حقيبة مجهّزة \\ \hline
$450$ & لعبة حاسوب \\ \hline
\else
$8$ & \\ \hline
$80$ & \\ \hline
$250$ & \\ \hline
$450$ & \\ \hline
\fi
\end{tabular}
\par} %
\filbreak

\item أمامكم قائمة هدايا ممكنة وقائمة مبالغ شراء تخوّل الحصول على هدايا. لائموا بين كل هدية ومبلغ الشراء الملائم، واكتبوا الإجابات في الجدول:

{\centering
\renewcommand{\arraystretch}{1.2}
\begin{tabular}{p{5cm}p{5cm}}
$1$. دفتر يوميات & \textenglish{A}. أكبر من $600$ \\
$2$. حقيبة مجهّزة & \textenglish{B}. $250 - 399$ \\
$3$. لعبة حاسوب & \textenglish{C}. $400 - 600$ \\
$4$. قسيمة شراء & \textenglish{D}. $10 - 100$ \\
\end{tabular}
\par} %

\vspace{1em}

{\centering
\renewcommand{\arraystretch}{1.4}
\begin{tabular}{|C{5cm}|C{4cm}|}
\hline
\textbf{الهدية التي يستحقها الزبون} & \textbf{مبلغ الشراء} \\
\hline
\ifwithsols
$1$. دفتر يوميات & \textenglish{D}. $10 - 100$ \\ \hline
$2$. حقيبة مجهّزة & \textenglish{B}. $250 - 399$ \\ \hline
$3$. لعبة حاسوب & \textenglish{C}. $400 - 600$ \\ \hline
$4$. قسيمة شراء & \textenglish{A}. أكبر من $600$ \\ \hline
\else
 & \\ \hline
 & \\ \hline
 & \\ \hline
 & \\ \hline
\fi
\end{tabular}
\par} %
\filbreak
\end{enumerate}
\filbreak

\end{enumerate}

\clearpage
\section*{القسم الثاني: ورقة عمل - نتدرّب على حل المشكلات}

\begin{enumerate}[resume, label=\textbf{تمرين \arabic*:}, wide=0pt, leftmargin=*, itemsep=2.2em]

% ---------------------------------------------------------------
\item \textbf{تمرين متابعة:} اكتبوا ما هو مُخرج مقطع البرنامج التالي دون تشغيله في برمجية بايثون. يجب الاستعانة بجدول متابعة.

\begin{boxCode}
\begin{english}
\begin{minted}[linenos]{python}
num = 1234
num1 = num // 100
num2 = num % 100
print("num1 = ", num1, " num2 = ", num2)
\end{minted}
\end{english}
\end{boxCode}

\ifwithsols
\begin{boxSolution}
\begin{center}
\begin{english}
\renewcommand{\arraystretch}{1.4}
\begin{tabular}{|C{1.2cm}|C{1.6cm}|C{1.6cm}|C{1.6cm}|C{5.2cm}|}
\hline
\textarabic{\textbf{السطر}} & \textbf{num} & \textbf{num1} & \textbf{num2} & \textarabic{\textbf{المخرجات}} \\
\hline
1 & 1234 & & & \\ \hline
2 & 1234 & 12 & & \\ \hline
3 & 1234 & 12 & 34 & \\ \hline
4 & 1234 & 12 & 34 & \texttt{num1 =\ \ 12\ \ num2 =\ \ 34} \\ \hline
\end{tabular}
\end{english}
\end{center}
\end{boxSolution}
\else
\begin{center}
\begin{english}
\renewcommand{\arraystretch}{1.4}
\begin{tabular}{|C{1.2cm}|C{1.6cm}|C{1.6cm}|C{1.6cm}|C{5.2cm}|}
\hline
\textarabic{\textbf{السطر}} & \textbf{num} & \textbf{num1} & \textbf{num2} & \textarabic{\textbf{المخرجات}} \\
\hline
 & & & & \\ \hline
 & & & & \\ \hline
 & & & & \\ \hline
 & & & & \\ \hline
\end{tabular}
\end{english}
\end{center}
\fi
\filbreak

% ---------------------------------------------------------------
\item \textbf{تمرين متابعة:} اكتبوا ما هو مُخرج مقطع البرنامج التالي دون تشغيله في برمجية بايثون. يجب الاستعانة بجدول متابعة.

\begin{boxCode}
\begin{english}
\begin{minted}[linenos]{python}
number = 123
print(number // 100)
print(number % 100 % 10)
print(number % 10)
\end{minted}
\end{english}
\end{boxCode}

\ifwithsols
\begin{boxSolution}
\begin{center}
\begin{english}
\renewcommand{\arraystretch}{1.4}
\begin{tabular}{|C{1.5cm}|C{2.5cm}|C{4cm}|}
\hline
\textarabic{\textbf{السطر}} & \textbf{number} & \textarabic{\textbf{المخرجات}} \\
\hline
1 & 123 & \\ \hline
2 & 123 & \texttt{1} \\ \hline
3 & 123 & \texttt{3} \\ \hline
4 & 123 & \texttt{3} \\ \hline
\end{tabular}
\end{english}
\end{center}
\end{boxSolution}
\else
\begin{center}
\begin{english}
\renewcommand{\arraystretch}{1.4}
\begin{tabular}{|C{1.5cm}|C{2.5cm}|C{4cm}|}
\hline
\textarabic{\textbf{السطر}} & \textbf{number} & \textarabic{\textbf{المخرجات}} \\
\hline
 & & \\ \hline
 & & \\ \hline
 & & \\ \hline
 & & \\ \hline
\end{tabular}
\end{english}
\end{center}
\fi
\filbreak

% ---------------------------------------------------------------
\item \textbf{تمرين متابعة:} اكتبوا ما هو مُخرج مقطع البرنامج التالي دون تشغيله في برمجية بايثون. يجب الاستعانة بجدول متابعة.

\begin{boxCode}
\begin{english}
\begin{minted}[linenos]{python}
num1 = 2
num2 = 3
num3 = 5
sum = num1 + num2 + num3
print(sum)
\end{minted}
\end{english}
\end{boxCode}

\ifwithsols
\begin{boxSolution}
\begin{center}
\begin{english}
\renewcommand{\arraystretch}{1.4}
\begin{tabular}{|C{1.5cm}|C{1.6cm}|C{1.6cm}|C{1.6cm}|C{1.6cm}|C{2.5cm}|}
\hline
\textarabic{\textbf{السطر}} & \textbf{num1} & \textbf{num2} & \textbf{num3} & \textbf{sum} & \textarabic{\textbf{المخرجات}} \\
\hline
1 & 2 & & & & \\ \hline
2 & 2 & 3 & & & \\ \hline
3 & 2 & 3 & 5 & & \\ \hline
4 & 2 & 3 & 5 & 10 & \\ \hline
5 & 2 & 3 & 5 & 10 & \texttt{10} \\ \hline
\end{tabular}
\end{english}
\end{center}
\end{boxSolution}
\else
\begin{center}
\begin{english}
\renewcommand{\arraystretch}{1.4}
\begin{tabular}{|C{1.5cm}|C{1.6cm}|C{1.6cm}|C{1.6cm}|C{1.6cm}|C{2.5cm}|}
\hline
\textarabic{\textbf{السطر}} & \textbf{num1} & \textbf{num2} & \textbf{num3} & \textbf{sum} & \textarabic{\textbf{المخرجات}} \\
\hline
 & & & & & \\ \hline
 & & & & & \\ \hline
 & & & & & \\ \hline
 & & & & & \\ \hline
 & & & & & \\ \hline
\end{tabular}
\end{english}
\end{center}
\fi
\filbreak

% ---------------------------------------------------------------
\item معطاة ثلاثة متغيرات \textenglish{a}، \textenglish{b}، \textenglish{c} وقيمها كالتالي:
\begin{center}
\textenglish{a = 10, b = 15, c = 0}
\end{center}
احسبوا قيمة كل من التعابير المنطقية في الجدول، واكتبوا القيمة (\textenglish{True} / \textenglish{False}) في العمود الفارغ.

{\centering
\renewcommand{\arraystretch}{1.4}
\begin{tabular}{|C{1.2cm}|>{\raggedright\arraybackslash}p{9.6cm}|C{3cm}|}
\hline
\textbf{البند} & \centering\textbf{التعبير المنطقي} & \textbf{قيمة التعبير المنطقي} \\
\hline
$1$ & \textenglish{\texttt{a >= b or b == c}} & \ifwithsols \textenglish{False} \fi \\ \hline
$2$ & \textenglish{\texttt{a \% 5 == 0 and b // 5 > c}} & \ifwithsols \textenglish{True} \fi \\ \hline
$3$ & \textenglish{\texttt{a < 0 or (b == c + 15 and a >= c)}} & \ifwithsols \textenglish{True} \fi \\ \hline
$4$ & \textenglish{\texttt{a \% 2 == c \% 2 and b \% 3 > a \% 3}} & \ifwithsols \textenglish{False} \fi \\ \hline
$5$ & \textenglish{\texttt{a != b and b != c}} & \ifwithsols \textenglish{True} \fi \\ \hline
\end{tabular}
\par} %

\ifwithsols
\begin{boxSolution}
\begin{enumerate}[label=\arabic*.]
\item $10 \geq 15$ خطأ، و $15 = 0$ خطأ، إذن \textenglish{False}.
\item $10 \% 5 = 0$ صحيح، و $15 // 5 = 3 > 0$ صحيح، إذن \textenglish{True}.
\item $10 < 0$ خطأ، لكن $15 = 0 + 15$ صحيح و $10 \geq 0$ صحيح، إذن \textenglish{True}.
\item $10 \% 2 = 0$ و $0 \% 2 = 0$ متساويان (صحيح)، لكن $15 \% 3 = 0$ ليس أكبر من $10 \% 3 = 1$ (خطأ)، إذن \textenglish{False}.
\item $10 \neq 15$ صحيح، و $15 \neq 0$ صحيح، إذن \textenglish{True}.
\end{enumerate}
\end{boxSolution}
\fi
\filbreak

% ---------------------------------------------------------------
\item في حانوت رياضيّ، في إطار حملة تُمنح هدايا للزبائن وفقًا لمبلغ الشراء.
يمكن أن تكون الهدية واحدة من الإمكانيات التالية: كرة يوغا (\textenglish{Yoga ball})، كرات تنس (\textenglish{Tennis balls})، كرة قدم (\textenglish{Football})، حذاء رياضي (\textenglish{Sports shoes}).
عند الشراء بمبلغ أقل من $80$ شيكل \textbf{لا يحصل} الزبون على هدية.

يتم تحديد الهدية حسب السيناريو الآتي (المتغير \textenglish{amount} هو مبلغ الشراء):

\begin{boxCode}
\begin{english}
\begin{minted}{python}
if amount < 80:
    print("Not eligible for a gift")
elif amount < 800:
    if amount >= 80 and amount <= 500:
        print("Yoga ball")
    elif amount < 650:
        print("Tennis balls")
    else:
        print("Football")
else:
    print("Sports shoes")
\end{minted}
\end{english}
\end{boxCode}

\begin{enumerate}[label=\alph*.]
\item قرّرت دانة تشغيل السيناريو. ساعدوها وأكملوا الجدول: لكل مبلغ، اكتبوا ما هي الهدية التي يستحقها الزبون.

{\centering
\renewcommand{\arraystretch}{1.4}
\begin{tabular}{|C{3cm}|C{6cm}|}
\hline
\textbf{مبلغ الشراء} & \textbf{الهدية التي يستحقها الزبون} \\
\hline
\ifwithsols
$25$ & لا يستحق هدية \\ \hline
$200$ & كرة يوغا \\ \hline
$750$ & كرة قدم \\ \hline
$999$ & حذاء رياضي \\ \hline
$550$ & كرات تنس \\ \hline
\else
$25$ & \\ \hline
$200$ & \\ \hline
$750$ & \\ \hline
$999$ & \\ \hline
$550$ & \\ \hline
\fi
\end{tabular}
\par} %
\filbreak

\item أمامكم قائمة الهدايا المُتاحة وقائمة مبالغ الشراء التي تؤهّل الحصول على هدية. لائموا لكل هدية مبلغ الشراء، واكتبوا الإجابات في الجدول التالي:

{\centering
\renewcommand{\arraystretch}{1.2}
\begin{tabular}{p{5cm}p{5cm}}
$1$. كرة يوغا & \textenglish{A}. أكبر من $800$ \\
$2$. كرات تنس & \textenglish{B}. $500 - 650$ \\
$3$. كرة قدم & \textenglish{C}. $650 - 800$ \\
$4$. حذاء رياضي & \textenglish{D}. $80 - 500$ \\
\end{tabular}
\par} %

\vspace{1em}

{\centering
\renewcommand{\arraystretch}{1.4}
\begin{tabular}{|C{5cm}|C{4cm}|}
\hline
\textbf{الهدية التي يستحقها الزبون} & \textbf{مبلغ الشراء} \\
\hline
\ifwithsols
$1$. كرة يوغا & \textenglish{D}. $80 - 500$ \\ \hline
$2$. كرات تنس & \textenglish{B}. $500 - 650$ \\ \hline
$3$. كرة قدم & \textenglish{C}. $650 - 800$ \\ \hline
$4$. حذاء رياضي & \textenglish{A}. أكبر من $800$ \\ \hline
\else
 & \\ \hline
 & \\ \hline
 & \\ \hline
 & \\ \hline
\fi
\end{tabular}
\par} %
\filbreak
\end{enumerate}
\filbreak

\end{enumerate}

\vspace{1cm}
\begin{flushleft}
أرجو لكم وقتًا ممتعًا.

الأستاذ محمود اغبارية.
\end{flushleft}

\end{document}
